% Shared paper content, included by main.tex (arXiv, named) and acl.tex (ARR review, anonymous). Edit here.

\begin{abstract}
Machine-translation evaluation uses LLMs to produce translations and, increasingly, to judge them; batching can change either model's response to the same request. On WMT24++, serving conditions changed \MtvBetweenMin{}--\MtvBetweenMax{} of segment translations, with more change in lower-resource directions. Metric-based system comparisons were mostly stable (\MtvSysCond{} of \MtvSysN{} changed), but a served LLM judge changed its error annotations on \MtvJudgeBetween{} of items. The preregistered per-comparison analysis flagged \MtvJudgeCond{} of \MtvJudgeN{} judge-based comparisons; a global Holm sensitivity leaves \MtvJudgeHolmCond{} nonsignificant score-order reversals and \MtvJudgeHolmSig{} significance changes. A parallel math-reasoning study found per-example answer changes but no condition-sensitive model comparisons. vLLM's batch-invariant mode produced identical token sequences across the tested conditions, at a model-dependent throughput cost. We measure that cost and evaluate a generic per-batch-size configuration as a practical mitigation for untuned models.
\end{abstract}

% OUTLINE
% 1 Introduction: output reproducibility in NLP; serving cost; generic-default result and scope.
% 2 Background: batch invariance; vLLM's mode on Hopper/Blackwell; the tuned table; #35404.
% 3 Setup and protocol: hardware, software, determinism and cost measurements, preregistered rules.
% 4 Where the cost comes from: tuned vs untuned stock cost; default tile config and compiled lookup.
% 5 Route 1, tuning the Triton kernel: pilot, serving benchmark, upstream PR.
% 6 Route 2, one generic default for untuned shapes.
% 7 Route 3, routing to cuBLAS: bitwise test, shape map, H200, serving determinism and cost; why not recommended.
% 8 Route 4, cuBLAS's batch-invariance mode on Blackwell.
% 9 A serving-level failure: Gemma-2 sliding window, diagnosis and fix.
% 10 Recommendations.
% 11 Limitations, related work, conclusion.

\section{Introduction}
MT evaluation uses LLM inference to generate translations and to assess their quality. Both stages are commonly batched, yet a serving system can change a response to the same prompt when unrelated requests share its batch or a prefix is cached: floating-point reduction order in matrix multiplication and attention can depend on batch shape \citep{he2025nondeterminism,yuan2025numerical}. This raises a field-level question: can ordinary serving conditions change MT evaluation records or the system-comparison verdicts drawn from them?

Our primary result is that they can change LLM-judge annotations and pairwise score ordering. In a pre-registered WMT24++ study \citep{deutsch2025wmt24pp}, between conditions, \MtvBetweenMin{}--\MtvBetweenMax{} of segment translations changed, with more change in lower-resource than high-resource directions (\MtvLowBetween{} vs. \MtvHighBetween{}). Metric-based comparisons were mostly stable, but the judge's annotations changed on \MtvJudgeBetween{} of items. The unadjusted per-comparison rule flagged \MtvJudgeCond{} of \MtvJudgeN{} judge-based system comparisons: \MtvJudgeSig{} changed significance and \MtvJudgeSign{} reversed the score order of a nonsignificant pair while identical repeats agreed. A post hoc Holm adjustment across all judge comparisons and runs leaves \MtvJudgeHolmCond{} condition-sensitive pairs, all nonsignificant score-order reversals; \MtvJudgeHolmSig{} significance changes remain. A matched math-reasoning study provides a contrast: \EvFlipBetweenMin{}--\EvFlipBetweenMax{} of examples changed correctness, yet \EvChangedCond{} of \EvComparisons{} system comparisons were condition-sensitive. Thus, the evaluation effect depends on what is measured: in this MT sample, serving variation changed the judge's item-level annotations and pairwise score ordering, while neural-metric comparisons largely absorbed it.

We then evaluate reproducible serving as a mitigation and measure its cost. vLLM's batch-invariant mode produced identical generated-token sequences across the tested conditions, but its throughput cost depends on model shapes and kernel configurations. We trace much of that cost to the default tile configuration, evaluate a generic per-batch-size fallback for untuned BF16 shapes, and document where attention and vendor-library behavior limit the recommendation. Our contributions are:
\begin{itemize}
\item An MT evaluation study showing that serving conditions change translation records and LLM-judge annotations and can change the judge's pairwise score ordering beyond identical-repeat variation, while metric-based comparisons are mostly stable (Section~\ref{sec:mtvar}).
\item A matched math-reasoning study that bounds the result: serving conditions changed per-example correctness but not model-comparison verdicts in the tested setup (Section~\ref{sec:mathvar}).
\item An assessment of token-level batch-invariant serving as a reproducibility measure, its model-dependent cost, and a generic fallback that recovers much of the tuning gain for untuned shapes; we also document the attention and vendor-library boundaries (Appendix~\ref{app:serving}).
\end{itemize}
Every pass/fail rule was written before its run. Code, raw logs and decision rules are available with the project artifact.

\section{Background}
A batch-invariant kernel computes each row of its output with the same sequence of floating-point operations regardless of how many other rows are in the batch. Reductions whose split depends on batch size (split-K matrix multiplication, split-KV attention, and some normalization kernels) break this property \citep{he2025nondeterminism}. vLLM's batch-invariant mode, enabled with \texttt{VLLM\_BATCH\_INVARIANT=1}, pins attention to a single split, replaces reductions with fixed-order kernels, disables reduced-precision reductions, prefers cuBLASLt, and pins collective communication to fixed algorithms. Unquantized Linear layers run a persistent Triton matrix-multiplication kernel on every CUDA GPU, a choice made after gpt-oss failed batch invariance with cuBLAS on H100 \citep{vllm35404}. The kernel's configuration comes from a table keyed by GPU family and weight shape $(N, K)$: each entry fixes one reduction tile \texttt{BLOCK\_K} for the shape and chooses the remaining tile sizes per batch-size bucket. Because \texttt{BLOCK\_K} alone fixes the reduction order, any entry of this form is batch-invariant by construction. Shapes without an entry use a single default configuration for every batch size.

\section{Setup and protocol}
The MT and math evaluation studies and the core-PR serving comparison use runtime vLLM build 0.34.0rc1.dev352 with the source revision tested by PR~\#67376 (commit \texttt{beae62efc9}, based on \texttt{df8c0296}). Earlier kernel-tuning and vendor-library measurements use their own matched vLLM builds; source and runtime metadata are preserved with each result set. Experiments use bf16 on H100 unless noted; H200 and B200 appear only in the comparisons identified below. \emph{Evaluation conditions}: requests run with cold or warm prefix cache, alone or mixed with unrelated requests, and with chunked prefill; every condition has identical repeats as a noise baseline. For the math and MT studies, outputs are compared token by token and top-20 log probabilities are also recorded. Prompts are few-shot GSM8K and MATH-500 problems for math, and WMT24++ segments for MT. A control with batch-invariant mode off must diverge for a test to count. The preregistered primary analysis uses an unadjusted per-comparison significance threshold; we additionally report a post hoc global Holm correction across all judge comparisons, conditions and repeats. These analyses are descriptive, and condition sensitivity additionally requires agreement between identical repeats. \emph{Cost}: decode throughput loss against the mode off; the paired serving comparisons keep the workload and source build fixed. \emph{Bitwise kernel tests}: a row multiplied alone is compared with the same row inside batches of 1 to 8192 rows. Pass/fail rules were written before each run (repository \texttt{PROTOCOL.md}).

\section{Does serving change evaluation conclusions?}\label{sec:evalvar}
\subsection{Machine translation}\label{sec:mtvar}
Five instruction-tuned LLMs (Gemma-2-9B, Gemma-3-12B, Aya-Expanse-8B, Qwen2.5-7B and Llama-3.1-8B) translated WMT24++ \citep{deutsch2025wmt24pp} from English into German, Chinese, Japanese, Icelandic, Swahili, Gujarati and Zulu, using every segment not flagged as a bad source, with a one-sentence translation prompt and greedy decoding. They ran under the four serving conditions described below for math (serial, full batch, mixed with unrelated prompts, chunked prefill), each twice, with batch-invariant mode off and on, on H100. We scored translations with MetricX-24 \citep{juraska-etal-2024-metricx}, COMET-22 \citep{rei-etal-2022-comet} and CometKiwi-23 \citep{rei-etal-2023-scaling}; the two COMET models do not cover Zulu. Each unique translation was scored once at a fixed batch size, so identical translations receive identical scores. As an LLM judge, Qwen2.5-14B-Instruct annotated translation errors in the GEMBA-MQM style \citep{kocmi2023gemba,kocmi2023gembamqm} for a fixed sample of segments in four directions from every system, and the judge itself ran under the same conditions and repeats. A comparison between two systems on one direction, scored by one metric or by the judge, is condition-sensitive under the same rule as for math, with a paired sign-flip permutation test in each run, in the spirit of paired significance testing for MT \citep{koehn-2004-statistical}.

\begin{table*}[t]\centering\small
% Generated by src/make_macros.py from results/stats.json; each row is checked against its archived runs.
\begin{tabular}{lrrrr}
\toprule
Model & Changed, between & Changed, within & BI-on differing & BI-on time (C2) \\
\midrule
Gemma-2-9B & 30.4\% & 12.1\% & 0 & +50.4\% \\
Gemma-3-12B & 32.8\% & 6.8\% & 0 & +76.9\% \\
Aya-Expanse-8B & 19.3\% & 0.0\% & 0 & +21.3\% \\
Qwen2.5-7B & 49.3\% & 11.8\% & 0 & +40.5\% \\
Llama-3.1-8B & 34.4\% & 0.0\% & 0 & +19.0\% \\
\midrule
Judge (Qwen2.5-14B) & 51.1\% & 0.5\% & 0 & +36.6\% \\
\bottomrule
\end{tabular}

\caption{Serving conditions and MT evaluation records. Changed: mean share of segment translations (judge: annotations) that differ between runs of different conditions (between) or between the two repeats of a condition (within), with batch-invariant mode off, averaged over directions. BI-on differing: outputs not identical across all 8 runs with the mode on. BI-on time: generation time against the mode off in the full-batch condition.}\label{tab:mtvar}
\end{table*}

Between conditions, \MtvBetweenMin{}--\MtvBetweenMax{} of segment translations changed, depending on the system; repeats of one condition differed by at most \MtvWithinMax{}, and not at all for two systems. Changes were more frequent in lower-resource directions (\MtvLowBetween{} for Icelandic, Swahili, Gujarati and Zulu against \MtvHighBetween{} for German, Chinese and Japanese; a descriptive grouping, not a pre-registered test). Averaged over segments, metric scores absorbed most of this: \MtvSysCond{} of \MtvSysN{} metric-based comparisons was condition-sensitive, a change in significance, while \MtvSysRepeat{} changed between identical repeats, and the system ranking differed across runs in \MtvRankCells{} direction--metric cells. The judge was different. Its own annotations changed on \MtvJudgeBetween{} of items between conditions (\MtvJudgeWithin{} between repeats), and its scores on \MtvJudgeScoreChange{}. \MtvJudgeCond{} of \MtvJudgeN{} judge-based comparisons were condition-sensitive and \MtvJudgeRepeat{} changed between repeats: \MtvJudgeSig{} changed significance, and \MtvJudgeSign{} reversed which system scored higher on a difference that was not significant, each between serial serving and the batched conditions. With batch-invariant mode on, \MtvBiDiffer{} token sequences differed across conditions for any translator or the judge, at a generation-time cost of \MtvTimeMin{}--\MtvTimeMax{}.

The metrics themselves are not batch-invariant. Scoring one system's translations at batch size one and sixteen changed \MtvSideMxShare{} of MetricX-24 segment scores (by at most \MtvSideMxMax{} points) and \MtvSideCometShare{} of COMET-22 scores (at most \MtvSideCometMax{}). Our design avoids this by scoring each translation once, but evaluations that rescore systems in different batches inherit it.

This contrast matters for MT practice: stable system-level metric comparisons can coexist with different translation records, while a served judge adds another inference stage whose annotations and verdicts may depend on batching. The condition-sensitive judge comparisons passed our repeat-agreement rule, so identical-repeat variation did not explain them. This is evidence about evaluation reproducibility in our setup, not evidence that one serving condition produces better translations.

\subsection{Math reasoning}\label{sec:mathvar}
As a second domain, we served six instruction-tuned models (Qwen2.5-7B, Llama-3.1-8B, Gemma-2-9B, Mistral-7B-v0.3, OLMo-2-7B and Phi-4-mini) on GSM8K test and MATH-500, zero-shot, with greedy decoding and up to 1024 new tokens, under four serving conditions in vLLM: serial (one request at a time), full batch (all prompts submitted together), mixed load (each benchmark prompt interleaved with an unrelated HotpotQA prompt), and chunked prefill with a 512-token budget. Each condition ran twice in separate processes, so that differences between conditions can be set against differences between identical repeats, and the whole design ran with batch-invariant mode off and on, on H100. A response is correct if its last boxed answer matches the reference after a fixed normalization; truncated responses count as wrong. For each of the 15 model pairs on each benchmark, we compute the accuracy difference and an exact McNemar test in every run. A comparison is condition-sensitive if its sign or its significance at the 0.05 level differs between conditions while the two repeats of each condition agree. The design, grader and decision rule were fixed before the runs.

\begin{table*}[t]\centering\small
% Generated by src/make_macros.py from results/stats.json; each row is checked against its archived runs.
\begin{tabular}{lrrrrrr}
\toprule
Model & GSM8K acc. & MATH-500 acc. & Flips between & Flips within & BI-on differing & BI-on time (C2) \\
\midrule
Qwen2.5-7B & 91.7--92.1 & 69.6--71.6 & 1.4\%/4.8\% & 0.5\%/2.2\% & 0 & +32.6\% \\
Llama-3.1-8B & 81.9--82.7 & 41.4--42.0 & 2.9\%/6.3\% & 0.0\%/0.4\% & 0 & +20.1\% \\
Gemma-2-9B & 87.4--87.9 & 44.2--45.8 & 0.9\%/2.5\% & 0.4\%/1.1\% & 0 & +47.7\% \\
Mistral-7B & 45.1--45.9 & 11.6--13.2 & 6.3\%/3.9\% & 0.0\%/0.3\% & 0 & +19.1\% \\
OLMo-2-7B & 71.9--72.5 & 27.0--27.8 & 3.1\%/4.0\% & 0.0\%/0.1\% & 0 & +23.7\% \\
Phi-4-mini & 89.8--90.0 & 56.6--58.0 & 1.5\%/5.4\% & 0.0\%/0.2\% & 0 & +42.7\% \\
\bottomrule
\end{tabular}

\caption{Serving conditions and evaluation records. Accuracy ranges span the 8 runs with batch-invariant mode off; flips are the mean share of examples whose correctness differs between runs of different conditions (between) or between the two repeats of a condition (within), for GSM8K/MATH-500. BI-on differing: requests whose tokens are not identical across all 8 runs with the mode on. BI-on time: generation time against the mode off in the full-batch condition.}\label{tab:evalvar}
\end{table*}

Serving conditions changed per-example records but not conclusions (Table~\ref{tab:evalvar}). Between runs in different conditions, \EvFlipBetweenMin{}--\EvFlipBetweenMax{} of examples changed correctness on average for each model and benchmark; the largest rate between two identical repeats of one condition was \EvFlipWithinMax{}. Accuracy across the 8 runs of a model spanned at most \EvSpreadMax{}. \EvChangedCond{} of \EvComparisons{} comparisons were condition-sensitive: the verdict of \EvChangedRepeat{} comparison differed between runs, and it also differed between identical repeats. The six models ranked identically in every run on both benchmarks. With exact-match grading and accuracy gaps that are mostly wide, serving conditions changed individual records without changing model comparisons, as did metric-based comparisons in MT. With batch-invariant mode on, served with the generic configuration of Section~\ref{sec:generic}, \EvBiDiffer{} requests differed across the 8 runs of any model, at a generation-time cost of \EvTimeCTwoMin{}--\EvTimeCTwoMax{} in the full-batch condition and \EvTimeCOneMin{}--\EvTimeCOneMax{} when serving one request at a time.

\section{A practical control for reproducible evaluation}\label{sec:generic}
Batch-invariant serving is a reproducibility control, not a translation-quality intervention: it holds generated token sequences fixed across the tested serving conditions, at a model-dependent throughput cost. We therefore evaluate a generic per-batch-size configuration as a separate mitigation for untuned BF16 shapes. It uses training shapes that exclude both evaluation models, together with breakable CUDA graphs so the compiled forward retains the runtime batch-size lookup. Our initial generic arm did not use the intended per-batch settings because the compiled forward froze the lookup; we voided that run and repeated it with breakable graphs and a stock-plus-breakable control. The correction is recorded in the protocol.

The model-level throughput arms install the selected configuration through an experiment hook in vLLM; the hook applies the same configuration and graph setting implemented by the open core change. An H100 smoke from the core PR \citep{vllm67376} confirmed that the generic QKV configuration was selected, breakable graphs were active, and the target prompt produced identical tokens alone and in a mixed batch. We then served the core PR against its exact merge base on the same runtime (Table~\ref{tab:corepr}; online \texttt{vllm bench serve}, fresh server per arm, \CoreRounds{} rotated rounds). The PR raised Qwen2.5-7B concurrency-32 decode throughput by \CoreQsevenDec{} and single-request decode by \CoreQsevenDecOne{}; the gain holds on H200 and for Qwen2.5-32B at TP=2. Mistral-7B, whose shapes are mostly covered by tuned entries for a model with the same projections, is unchanged (\CoreMistralDec{} decode). Gemma-2-9B gains on decode but loses \CoreGemmaPreLoss{} on prefill (\CoreGemmaPreWins{} of \CoreRounds{} rounds), and its long-prompt decode does not improve. A separate diagnostic run (\DiagGemmaRounds{} rounds) isolates the cause as the graph mode, not the tiles. Forcing breakable graphs on the merge base, with no new tiles, slows Gemma-2-9B prefill by \DiagGemmaBasebrkPre{} and decode by \DiagGemmaBasebrkDec{} (\DiagGemmaBasebrkPreWins{} wins each). The PR with breakable graphs turned off is \DiagGemmaPrnobrkPre{} faster on prefill (\DiagGemmaPrnobrkPreWins{} of \DiagGemmaRounds{} rounds) but gains only \DiagGemmaPrnobrkDec{} on decode, against \DiagGemmaPrDec{} with them; the PR as submitted is \DiagGemmaPrPre{} on prefill in this run. The generic tiles therefore help Gemma-2-9B in both phases, and the prefill loss is the cost of running it without torch.compile, which breakable graphs require. Whether to enable them is then a decode-versus-prefill trade-off per model rather than a property of the configuration.

\begin{table*}[t]
\centering\small
% Generated by src/make_macros.py from results/stats.json; each row is checked against its raw vllm bench serve records.
\begin{tabular}{lrrrr}
\toprule
Model, GPU & Decode, c32 & Prefill, c32 & Decode, c1 & Long prompt, c32 \\
\midrule
Qwen2.5-7B, H100 & +93.4\% (3/3) & +4.4\% (3/3) & +153.4\% (3/3) & +26.0\% (3/3) \\
Qwen2.5-7B, H200 & +77.2\% (3/3) & +6.8\% (3/3) & +186.5\% (3/3) & +23.6\% (3/3) \\
Qwen2.5-32B, TP=2, H100 & +55.6\% (3/3) & +6.3\% (3/3) & +137.9\% (3/3) & +52.7\% (3/3) \\
Gemma-2-9B, H100 & +40.5\% (3/3) & $-$8.1\% (0/3) & +92.3\% (3/3) & $-$0.4\% (1/3) \\
Mistral-7B, H100 & +0.9\% (2/3) & +1.5\% (3/3) & +2.0\% (2/3) & $-$0.8\% (1/3) \\
\bottomrule
\end{tabular}

\caption{Core PR against its merge base, batch-invariant mode on: median of \CoreRounds{} paired per-round throughput ratios (output tokens for decode, requests for prefill), with PR wins. c32 and c1: maximum concurrency. Raw logs are in the artifact.}\label{tab:corepr}
\end{table*}

The held-out bitwise check found \GenericCheckBad{} differing rows among \GenericCheckRows{} across \GenericTestShapes{} shapes; in the serving crossover, \GenericCrossMismatch{} of \GenericCrossRequests{} requests changed tokens under the generic configuration. At batch 32, the generic default reduces Qwen2.5-7B's loss from \GenericStockB{} to \GenDefaultLossBatchThirtyTwo{}, close to the per-model tuned result of \GenTunedLossBatchThirtyTwo{} and recovering \GenericRecovered{} of the stock-to-tuned gap. Breakable graphs alone do not explain this gain: the stock-plus-breakable control has a loss of \GenericOnbreakB{}. On Gemma-2-9B, the generic default lowers loss from \GenericGemmaStockB{} to \GenericGemmaB{}; the stock-plus-breakable control is \GenericGemmaOnbreakB{}. At batch 128, the generic default's loss is \GenDefaultLossBatchOneTwentyEight{} compared with \GenTunedLossBatchOneTwentyEight{} for the per-model table. The detailed batch-size curves and per-model tuning comparison are in Appendix~\ref{app:serving}.

\section{Recommendations}\label{sec:rec}
\begin{table*}[t]\centering\small
\begin{tabular}{llll}
\toprule
Route & Batch-32 loss & Invariance & Status \\
\midrule
Stock, default config & \QsevenStock{} & by construction & vLLM default \\
Generic + breakable graphs & \GenDefaultLossBatchThirtyTwo{} & by construction & evaluated \\
Per-model tuned & \TunedB{} & by construction & measured ref. \\
cuBLAS routing & \QsevenPlain{} & shape-dependent & avoid \\
\bottomrule
\end{tabular}
\caption{Measured routes to batch-invariant serving for Qwen2.5-7B on H100.}\label{tab:rec}
\end{table*}
For practitioners who need reproducible outputs today (Table~\ref{tab:rec}): use vLLM's batch-invariant mode and check whether the model's Linear shapes are in the tuned table for the GPU. For untuned dense BF16 shapes on H100, the generic configuration with breakable graphs is a promising default; where maintained shape entries exist, use those. Do not route to cuBLAS on Hopper. On Blackwell, cuBLAS's batch-invariance mode is promising if vLLM avoids applying the redundant PyTorch reduced-precision setting when the cuBLAS mode is active; end-to-end throughput remains unmeasured (Appendix~\ref{app:blackwell}).

Implementation details and additional diagnostics are in Appendix~\ref{app:serving}. For maintainers of batch-invariant serving, three lessons from this study: test invariance per shape, including narrow layers such as routers and indexers, since wide projections passing says little about them; treat any invariance property of a vendor library as an assumption to be re-checked on each GPU generation and library release unless the vendor documents it; and when a library offers its own invariance mode, check how it interacts with the framework's other numerical settings, which can silently undo it.

\section{Related Work}
\citet{he2025nondeterminism} identify batch-size-dependent reductions as a source of inference nondeterminism; vLLM and SGLang address it with batch-invariant and deterministic modes \citep{sglang2025deterministic}. Detent reduces the cost of this goal by fixing batch-dependent values used in kernel selection \citep{xin2026detent}. Vosti formalizes a stronger, logit-level guarantee and reports that existing production modes can produce different logits under some execution variations \citep{qin2026vosti}; our empirical claims concern generated-token identity and evaluation outcomes, not logit equality. Other work studies numerical variation in reasoning accuracy \citep{yuan2025numerical}, verified speculation \citep{gond2026llm42}, and run-to-run evaluation variance \citep{madaan2024variance,biderman2024lessons,ouyang2023chatgpt,atil2024nondeterminism,song2024greedy}. MT evaluation uses neural metrics \citep{rei-etal-2022-comet,juraska-etal-2024-metricx,rei-etal-2023-scaling}, LLM judges \citep{kocmi2023gemba,kocmi2023gembamqm}, and significance tests over system scores \citep{koehn-2004-statistical,kocmi-etal-2021-ship}. We isolate serving conditions under a fixed engine, model and decoding configuration, compare their effects with identical repeats, and test whether they change MT judge-based system comparisons. Tuned configurations for vLLM's batch-invariant kernel have been contributed for individual models \citep{vllm62219,vllm65903}; our cost measurements assess these approaches as practical mitigations.

\section{Conclusion}
In this WMT24++ study, ordinary serving conditions changed translations and a served LLM judge's annotations. The preregistered per-comparison analysis flagged condition-sensitive judge comparisons; under global Holm correction, the remaining differences are nonsignificant score-order reversals, not changes to significant winners. Neural-metric comparisons were largely stable. Identical repeats let us distinguish condition effects from ordinary run-to-run variation. These are findings about MT evaluation reproducibility, not evidence that a serving condition improves translation quality. A matched math study provides a contrast, where per-example answers changed but model comparisons did not. Token-level batch-invariant serving is a practical way to control this source of variation in the tested setups, though its throughput cost depends on the model. The generic configuration reduces that cost for untuned shapes in our serving comparison; its design and limitations are documented in Appendix~\ref{app:serving}.

\paragraph{Future work: from reproducible evaluation to better translation.} Our results suggest four routes from reproducible serving to translation quality. First, serving noise sets a floor on the smallest improvement a comparison can detect; removing it at low cost lets work on tokenizers, prompts or fine-tuning resolve smaller real gains with fewer repeated runs. Second, quality-estimation reranking selects among candidate translations with a metric or an LLM judge, so if those scores depend on batching, the selection does too; deterministic scoring would make such gains repeatable. Third, back-translation and distillation for low-resource languages generate training data with served models, and batch-invariant serving would make that data exactly regenerable. Fourth, translations changed most in lower-resource directions, which suggests that decoding there is less stable; fixes for it can now be evaluated without serving noise. We leave these to future work.

\label{end-of-content}

\section{Limitations}
We tested one engine, single-node tensor parallelism up to four GPUs, dense models up to 32B parameters, and bf16; we did not test multi-node serving or production traffic. The hook-based evaluation of the generic default used two held-out models at tensor parallelism one on H100; the core-PR serving comparison adds H200, tensor parallelism two and a partly tuned model, but covers four dense model families. The Gemma-2-9B prefill regression is attributed to the graph mode by a single diagnostic run, and we did not test whether other models with the same architecture share it. The tuned-table decode curves in the main comparison use three repeats; the PR branch breakable-graph comparison has two. The branch result covers a quantized gpt-oss model whose tested unquantized projections use the table; it does not establish behavior for other quantized models. Sparse MLA is outside our tested support: current vLLM rejects GLM-5.3 under batch invariance because no sparse backend opts in, and a separate report describes batch-dependent outputs under measurement-only opt-in \citep{vllm65278}. cuBLAS behavior was measured on the library versions listed and may change. The MT study covers five open translators of similar size, seven WMT24++ directions from English, one judge model and one prompt per task; judges based on larger or proprietary models, and other prompts, may vary differently. The math study covers two benchmarks with exact-match grading, greedy decoding, six instruction-tuned models from Phi-4-mini to Gemma-2-9B, one GPU type and one engine; comparisons between models with closer accuracy, sampled decoding, or tasks scored by model judges may be more sensitive to serving conditions. We compare only against vLLM's own batch-invariant mode, not against other engines' deterministic modes or LLM-42 \citep{gond2026llm42}.

\section*{Use of AI Assistants}
An AI coding assistant (Claude Code) was used to write experiment and analysis code, run experiments, and draft text. The authors reviewed the code, the raw outputs behind every reported number, and the final text.
